\section{Algoritma Euclid dan Euclid Diperluas}

\subsection{Algoritma Euclid}
Algoritma Euclid menghitung $\gcd(a, b)$ secara efisien \cite{stinson2018,menezes1996}. Ide: $\gcd(a, b) = \gcd(b, a \bmod b)$. Proses berulang hingga $b = 0$, maka $\gcd = a$.

\subsection{Algoritma Euclid Diperluas}
Algoritma Euclid diperluas tidak hanya menghitung $\gcd(a, b)$ tetapi juga koefisien $x$ dan $y$ sehingga $ax + by = \gcd(a, b)$. Koefisien $x$ memberikan invers modular dari $a$ modulo $b$ jika $\gcd(a, b) = 1$ \cite{stinson2018,menezes1996}.

\begin{lstlisting}[style=cstyle,caption={Algoritma Euclid Diperluas dalam C}]
int extended_gcd(int a, int b, int *x, int *y) {
    if (b == 0) { *x = 1; *y = 0; return a; }
    int x1, y1;
    int g = extended_gcd(b, a % b, &x1, &y1);
    *x = y1;
    *y = x1 - (a / b) * y1;
    return g;
}
\end{lstlisting}

Implementasi lengkap dalam MATLAB, C++, dan Pascal tersedia di Bab 14.
